\chapter{Všeobecné pravidlá hry}

\section*{Cieľ hry}
Každá rola má vlastný cieľ:
\begin{itemize}
    \item \textbf{Šerif}: vyradiť všetkých banditov a odpadlíka.
    \item \textbf{Banditi}: vyradiť šerifa.
    \item \textbf{Pomocníci šerifa}: chrániť šerifa a pomôcť mu splniť jeho cieľ.
    \item \textbf{Odpadlík}: zostať poslednou postavou v hre a nakoniec vyradiť šerifa.
\end{itemize}

\section*{Príprava hry}
\begin{enumerate}
    \item Rozdeľte role podľa počtu hráčov a každému dajte jednu kartu role. Šerif svoju rolu odhalí, ostatné role ostávajú skryté.
    \item Každému hráčovi prideľte jednu postavu a počet životov podľa počtu guliek na karte postavy. Šerif začína s jedným životom navyše.
    \item Zamiešajte hrací balíček a každému hráčovi rozdeľte toľko kariet, koľko má počiatočných životov.
    \item Zvyšné karty vytvoria dobrací balíček. Vedľa neho nechajte miesto na odhadzovaciu kopu.
\end{enumerate}

\section*{Priebeh ťahu}
Hra prebieha po smere hodinových ručičiek a šerif začína. Každý ťah má tri fázy:
\begin{enumerate}
    \item \textbf{Potiahnuť 2 karty.}
    \item \textbf{Zahrať ľubovoľný počet kariet.}
    \item \textbf{Odhodiť prebytočné karty.}
\end{enumerate}

Na konci ťahu môže mať hráč na ruke najviac toľko kariet, koľko má práve životov.

\section*{Obmedzenia pri hraní kariet}
Počas svojho ťahu platia tieto základné obmedzenia:
\begin{itemize}
    \item možno zahrať iba 1 kartu \textbf{Bang!} za ťah; toto obmedzenie sa vzťahuje len na konkrétnu kartu Bang!, nie na všetky efekty Bang!
    \item pred sebou nemožno mať 2 karty s rovnakým názvom
    \item v hre možno mať iba 1 zbraň
\end{itemize}

Hnedé karty sa po zahraní odhodia. Modré a oranžové karty ostávajú vyložené v hre, pokiaľ nie sú odstránené alebo samy neopustia hru.

\section*{Vzdialenosť a dostrel}
Vzdialenosť medzi susednými hráčmi je 1. Medzi dvoma hráčmi sa vždy počíta kratšia cesta po stole. Keď je niektorý hráč vyradený, pri počítaní vzdialenosti sa už neberie do úvahy.

Zbrane menia iba dostrel hráča. Samy o sebe nemenia vzdialenosť medzi hráčmi. Vzdialenosť upravujú iba efekty, ktoré to výslovne hovoria, napr. Mustang alebo Scope.

\section*{Strata životov, vyradenie a odmeny}
Ak hráč stratí posledný život, je vyradený z hry, pokiaľ sa okamžite nezachráni kartou Pivo alebo iným pravidlom, ktoré to výslovne povoľuje.

Keď je hráč vyradený:
\begin{itemize}
    \item odhalí svoju rolu
    \item odhodí všetky karty z ruky aj z hry
    \item vyhodnotia sa prípadné odmeny alebo tresty za jeho vyradenie
\end{itemize}

Základné odmeny a tresty:
\begin{itemize}
    \item ak šerif vyradí pomocníka, musí odhodiť všetky svoje karty z ruky aj z hry
    \item hráč, ktorý vyradí banditu, si potiahne 3 karty z balíčka
\end{itemize}

\section*{Koniec hry}
Hra končí, keď nastane jedna z týchto podmienok:
\begin{itemize}
    \item ak je vyradený šerif a odpadlík nie je posledný živý hráč, vyhrávajú banditi
    \item ak sú vyradení všetci banditi aj odpadlík, vyhráva šerif a jeho pomocníci
    \item ak je šerif vyradený a odpadlík zostane posledný v hre, vyhráva odpadlík
\end{itemize}

\section*{Dôležité všeobecné pripomienky}
\begin{itemize}
    \item Kartu \textbf{Vedle!} a kartu \textbf{Pivo} možno hrať aj mimo vlastného ťahu, ale iba v prípadoch, ktoré dovoľujú základné pravidlá alebo iné výslovné pravidlo.
    \item Karta \textbf{Pivo} nemá účinok, ak sú v hre už len 2 hráči.
    \item \textbf{Saloon} nie je karta Pivo, preto sa neriadi automaticky rovnakými obmedzeniami.
    \item Počas duelu sa karty Bang! odhadzujú, nie hrajú, preto sa nepočítajú do limitu jednej karty Bang! za ťah.
    \item Keď pravidlo vyžaduje konkrétnu kartu Bang! alebo Vedle!, nejde automaticky o ľubovoľný podobný efekt, pokiaľ to pravidlo výslovne nehovorí.
    \item Kým sa úplne nedokončí jeden efekt karty alebo schopnosti, nemožno začať hrať ďalšiu kartu.
    \item Karty role a postavy nie sú počas hry cieľom kariet ani schopností, pokiaľ nejaké pravidlo výslovne nehovorí opak.
    \item Pri vyradení kartami ako Indiáni! alebo Guľomet sa za vyradenie považuje hráč, ktorý túto kartu zahral; pri Dynamite to neplatí.
\end{itemize}

\section*{Varianta Dynamite Box}
Ak hráte variant s dreveným dynamitom, platia tieto dodatky:
\begin{itemize}
    \item na konci prípravy dostane šerif drevený dynamit pred seba
    \item vždy, keď počas svojho ťahu zasiahnete iného hráča kartou Bang!, vezmete si drevený dynamit pred seba
    \item keď sa v balíčku objaví karta Dynamit, nepočíta sa ako normálne potiahnutá alebo ťahaná karta; namiesto nej sa ukáže a položí pred hráča, ktorý má drevený dynamit
    \item od tej chvíle sa karta Dynamit správa podľa bežných pravidiel
    \item ak karta Dynamit opustí hru výbuchom alebo odhodením, hneď sa zamieša späť do balíčka
\end{itemize}
